% written by analysis/copy_checks_numbers.py from benchmarks/copy_checks
\newcommand{\MonoCultSSUAll}{17}
\newcommand{\MonoSpSSUAll}{28}
\newcommand{\MonoCultSSUMain}{19}
\newcommand{\MonoSpSSUMain}{28}
\newcommand{\MonoCultITSAll}{18}
\newcommand{\MonoSpITSAll}{24}
\newcommand{\MonoCultITSMain}{24}
\newcommand{\MonoSpITSMain}{29}
\newcommand{\MonoCultLSUAll}{19}
\newcommand{\MonoSpLSUAll}{28}
\newcommand{\MonoCultLSUMain}{22}
\newcommand{\MonoSpLSUMain}{30}
\newcommand{\MonoCultLinkedAll}{23}
\newcommand{\MonoSpLinkedAll}{31}
\newcommand{\MonoCultLinkedMain}{28}
\newcommand{\MonoSpLinkedMain}{35}
% cultures with >= 2 copies
\newcommand{\MonoCultN}{139}
\newcommand{\MonoCultPctMin}{12\%}
\newcommand{\MonoCultPctMax}{20\%}
% median share of a non-monophyletic culture's copies outside its largest pure clade (linked, main type)
\newcommand{\MonoCultEscMedian}{67\%}
% species monophyletic on the linked tree only once divergent copies are pruned
\newcommand{\MonoFixedN}{4}
\newcommand{\MonoComplexLinkedMain}{4/4}
\newcommand{\ArtShortParalog}{0}
\newcommand{\ArtShortRecDel}{22}
\newcommand{\ArtShortUnres}{0}
\newcommand{\ArtShortArtefact}{21}
\newcommand{\ArtShortArtefactTen}{2}
\newcommand{\ArtSameParalog}{3}
\newcommand{\ArtSameRecDel}{0}
\newcommand{\ArtSameUnres}{2}
\newcommand{\ArtSameArtefact}{3}
\newcommand{\ArtSameArtefactTen}{0}
% long-deletion threshold, nt
\newcommand{\ArtLConserved}{3}
\newcommand{\ArtLSpacer}{5}
\newcommand{\ArtMainFlagged}{5}
% post hoc: short-5.8S copies whose 5.8S deletion (same length, start within 2 nt) occurs under another label
\newcommand{\ShDelRecur}{43}
\newcommand{\ShDelRecurArtefact}{21}
\newcommand{\ShSixCopies}{40}
\newcommand{\ShSixMain}{1}
\newcommand{\ShSixCultures}{23}
% function check, ssu_flank, short-5.8S copies
\newcommand{\FnSSUIrr}{7}
\newcommand{\FnSSUSubs}{0.07}
\newcommand{\FnSSUSubsBase}{0.00}
\newcommand{\FnSSUDeg}{3}
\newcommand{\FnSSUBrokenShort}{10.0\%}
\newcommand{\FnSSUBrokenBase}{10.6\%}
\newcommand{\FnSSUDiff}{+0.0}
\newcommand{\FnSSUP}{0.17}
% function check, s58, short-5.8S copies
\newcommand{\FnFiveIrr}{25}
\newcommand{\FnFiveSubs}{2.00}
\newcommand{\FnFiveSubsBase}{0.00}
\newcommand{\FnFiveDeg}{2}
\newcommand{\FnFiveBrokenShort}{16.0\%}
\newcommand{\FnFiveBrokenBase}{4.0\%}
\newcommand{\FnFiveDiff}{+12.0}
\newcommand{\FnFiveP}{$<10^{-8}$}
% function check, lsu_flank, short-5.8S copies
\newcommand{\FnLSUIrr}{25}
\newcommand{\FnLSUSubs}{4.62}
\newcommand{\FnLSUSubsBase}{0.12}
\newcommand{\FnLSUDeg}{12}
\newcommand{\FnLSUBrokenShort}{7.6\%}
\newcommand{\FnLSUBrokenBase}{8.5\%}
\newcommand{\FnLSUDiff}{+0.8}
\newcommand{\FnLSUP}{0.07}
\newcommand{\FnLSUDegPct}{28\%}
\newcommand{\FnShortN}{43}
\newcommand{\FnSameSSUDeg}{0}
\newcommand{\FnSameFiveDeg}{0}
\newcommand{\FnSameLSUDeg}{0}
\newcommand{\FnVerdictShort}{structurally degraded}
\newcommand{\FnVerdictSame}{structurally intact}
% short copies whose 5.8S deletion removes paired positions of the 121-130 hairpin
\newcommand{\FnHairpinShort}{40}
\newcommand{\FnHairpinMain}{2}
% short copies whose 5.8S deletion removes only unpaired positions
\newcommand{\FnUnpairedShort}{3}
% loops of the shortened 5.8S hairpin in the short copies (post hoc)
\newcommand{\FnLoopUUCG}{25}
\newcommand{\FnLoopUUUG}{10}
\newcommand{\FnLoopCUCG}{5}
\newcommand{\FnLoopOther}{0}
% short copies whose shortened 6-bp stem is all Watson-Crick
\newcommand{\FnStemAllWC}{36}
\newcommand{\FnITSTwoBPShort}{0.16}
\newcommand{\FnITSTwoBPBase}{0.00}
\newcommand{\FnITSTwoMFEShort}{-0.23}
\newcommand{\FnITSTwoMFEBase}{-0.26}
% ssu_flank: share of main-type copies with a congener within the within-species spread
\newcommand{\ThrSSUCongener}{80.1\%}
\newcommand{\ThrSSUSameCulture}{0.00\%}
\newcommand{\ThrSSUSameSpecies}{0.00\%}
\newcommand{\ThrSSUOtherGenus}{2.7\%}
\newcommand{\ThrSSUOtherGenusSpread}{0.0\%}
% full_its: share of main-type copies with a congener within the within-species spread
\newcommand{\ThrITSCongener}{39.5\%}
\newcommand{\ThrITSSameCulture}{0.41\%}
\newcommand{\ThrITSSameSpecies}{0.66\%}
\newcommand{\ThrITSOtherGenus}{21.8\%}
\newcommand{\ThrITSOtherGenusSpread}{0.0\%}
% LSU_LROR-FLR2: share of main-type copies with a congener within the within-species spread
\newcommand{\ThrLSUCongener}{54.5\%}
\newcommand{\ThrLSUSameCulture}{0.14\%}
\newcommand{\ThrLSUSameSpecies}{0.14\%}
\newcommand{\ThrLSUOtherGenus}{11.1\%}
\newcommand{\ThrLSUOtherGenusSpread}{0.0\%}
% SSU_VT_NS31-AM1: share of main-type copies with a congener within the within-species spread
\newcommand{\ThrVTCongener}{93.7\%}
\newcommand{\ThrVTSameCulture}{0.00\%}
\newcommand{\ThrVTSameSpecies}{0.00\%}
\newcommand{\ThrVTOtherGenus}{3.9\%}
\newcommand{\ThrVTOtherGenusSpread}{13.5\%}
% within-species spread on the VT fragment (99th pct)
\newcommand{\ThrVTSpread}{1.7\%}
% VT fragment: share of main-type copies with a congener / another genus within 3%
\newcommand{\ThrVTCongenerDB}{98.7\%}
\newcommand{\ThrVTOtherGenusDB}{15.2\%}
\newcommand{\ThrVTOtherGenusN}{121}
\newcommand{\ThrVTOtherGenusDen}{794}
% of them, copies outside Cetraspora, Gigaspora and Dentiscutata
\newcommand{\ThrVTOtherGenusNotGiga}{4}
% full ITS: share of main-type copies more than 1.5% from every other main-type copy of their culture
\newcommand{\ThrITSSameCultureAboveSH}{24.4\%}
\newcommand{\ThrITSSameSpeciesAboveSH}{34.4\%}
\newcommand{\ThrITSCongenerDB}{1.8\%}
\newcommand{\ThrPairsN}{5{,}196}
% culture pairs with SSU distance above half the ITS distance
\newcommand{\ThrPairsAbove}{6}
\newcommand{\ThrPairsAboveMaxSSU}{0.48\%}
\newcommand{\ThrPairsRuleA}{0}
% median SSU-flank / full-ITS distance, culture pairs >= 2% apart in ITS
\newcommand{\ThrPairsRatio}{0.14}
% clades made only of short-5.8S copies in the Rhizophagus clade (SH-aLRT/UFBoot)
\newcommand{\TreeShortClades}{\spn{Rhizophagus irregularis} with \spn{{[}Rhizoglomus{]} vesiculiferum} (26 copies and 1 copy filed as main type; 100/92); \spn{Rhizophagus intraradices} with \spn{Rhizophagus prolifer} (8 copies; 100/99); \spn{Rhizophagus silesianus} (3 copies; 100/100); \spn{Rhizophagus irregularis} with \spn{Rhizophagus} sp. (2 copies; 100/100); \spn{Rhizophagus irregularis} (1 copy)}
\newcommand{\TreeShortCladesN}{5}
% short copies in the Rhizophagus clade; of them, in short-only clades of two or more tips; in such clades of two taxa
\newcommand{\TreeShortTotal}{40}
\newcommand{\TreeShortInMulti}{39}
\newcommand{\TreeShortMultiN}{4}
\newcommand{\TreeShortCross}{36}
\newcommand{\TreeShortCrossN}{3}
\newcommand{\TreeRhizTips}{103}
\newcommand{\TreeRhizCopies}{154}
% Kruger fragment (SSUmAf1 site to LSUmBr) and whole copy: lengths, nt
\newcommand{\LongKrMedian}{1{,}577}
\newcommand{\LongKrMin}{1{,}451}
\newcommand{\LongKrMax}{1{,}678}
\newcommand{\LongFullMedian}{2{,}758}
% within-species spread (99th pct)
\newcommand{\LongKrSpread}{3.3\%}
\newcommand{\LongFullSpread}{1.9\%}
% full_its at 97% identity: identity, split from culture, split from species, joined (all; outside complexes)
\newcommand{\LongITSNSId}{97.00\%}
\newcommand{\LongITSNSSplit}{10.6\%}
\newcommand{\LongITSNSSplitSp}{3.4\%}
\newcommand{\LongITSNSJoin}{6.5\%}
\newcommand{\LongITSNSJoinX}{0.0\%}
% full_its at 98.5% identity: identity, split from culture, split from species, joined (all; outside complexes)
\newcommand{\LongITSNEId}{98.50\%}
\newcommand{\LongITSNESplit}{24.4\%}
\newcommand{\LongITSNESplitSp}{13.7\%}
\newcommand{\LongITSNEJoin}{1.8\%}
\newcommand{\LongITSNEJoinX}{0.0\%}
% full_its at best single threshold: identity, split from culture, split from species, joined (all; outside complexes)
\newcommand{\LongITSBestId}{96.55\%}
\newcommand{\LongITSBestSplit}{8.7\%}
\newcommand{\LongITSBestSplitSp}{2.1\%}
\newcommand{\LongITSBestJoin}{7.8\%}
\newcommand{\LongITSBestJoinX}{0.0\%}
% full_its at best single threshold, complexes aside: identity, split from culture, split from species, joined (all; outside complexes)
\newcommand{\LongITSBestXId}{94.95\%}
\newcommand{\LongITSBestXSplit}{3.6\%}
\newcommand{\LongITSBestXSplitSp}{0.7\%}
\newcommand{\LongITSBestXJoin}{32.4\%}
\newcommand{\LongITSBestXJoinX}{3.1\%}
% full_its at within-species spread: identity, split from culture, split from species, joined (all; outside complexes)
\newcommand{\LongITSSprId}{94.01\%}
\newcommand{\LongITSSprSplit}{2.7\%}
\newcommand{\LongITSSprSplitSp}{0.5\%}
\newcommand{\LongITSSprJoin}{39.5\%}
\newcommand{\LongITSSprJoinX}{11.8\%}
% kruger at 97% identity: identity, split from culture, split from species, joined (all; outside complexes)
\newcommand{\LongKrNSId}{97.00\%}
\newcommand{\LongKrNSSplit}{4.1\%}
\newcommand{\LongKrNSSplitSp}{0.6\%}
\newcommand{\LongKrNSJoin}{19.1\%}
\newcommand{\LongKrNSJoinX}{0.3\%}
% kruger at 98.5% identity: identity, split from culture, split from species, joined (all; outside complexes)
\newcommand{\LongKrNEId}{98.50\%}
\newcommand{\LongKrNESplit}{9.3\%}
\newcommand{\LongKrNESplitSp}{2.9\%}
\newcommand{\LongKrNEJoin}{1.1\%}
\newcommand{\LongKrNEJoinX}{0.0\%}
% kruger at best single threshold: identity, split from culture, split from species, joined (all; outside complexes)
\newcommand{\LongKrBestId}{98.15\%}
\newcommand{\LongKrBestSplit}{6.7\%}
\newcommand{\LongKrBestSplitSp}{1.7\%}
\newcommand{\LongKrBestJoin}{7.0\%}
\newcommand{\LongKrBestJoinX}{0.0\%}
% kruger at best single threshold, complexes aside: identity, split from culture, split from species, joined (all; outside complexes)
\newcommand{\LongKrBestXId}{96.55\%}
\newcommand{\LongKrBestXSplit}{3.0\%}
\newcommand{\LongKrBestXSplitSp}{0.5\%}
\newcommand{\LongKrBestXJoin}{32.1\%}
\newcommand{\LongKrBestXJoinX}{2.6\%}
% kruger at within-species spread: identity, split from culture, split from species, joined (all; outside complexes)
\newcommand{\LongKrSprId}{96.70\%}
\newcommand{\LongKrSprSplit}{3.3\%}
\newcommand{\LongKrSprSplitSp}{0.6\%}
\newcommand{\LongKrSprJoin}{25.6\%}
\newcommand{\LongKrSprJoinX}{2.1\%}
% full at 97% identity: identity, split from culture, split from species, joined (all; outside complexes)
\newcommand{\LongFullNSId}{97.00\%}
\newcommand{\LongFullNSSplit}{1.1\%}
\newcommand{\LongFullNSSplitSp}{0.1\%}
\newcommand{\LongFullNSJoin}{77.1\%}
\newcommand{\LongFullNSJoinX}{36.0\%}
% full at 98.5% identity: identity, split from culture, split from species, joined (all; outside complexes)
\newcommand{\LongFullNEId}{98.50\%}
\newcommand{\LongFullNESplit}{5.1\%}
\newcommand{\LongFullNESplitSp}{1.2\%}
\newcommand{\LongFullNEJoin}{15.6\%}
\newcommand{\LongFullNEJoinX}{0.0\%}
% full at best single threshold: identity, split from culture, split from species, joined (all; outside complexes)
\newcommand{\LongFullBestId}{98.90\%}
\newcommand{\LongFullBestSplit}{7.0\%}
\newcommand{\LongFullBestSplitSp}{2.0\%}
\newcommand{\LongFullBestJoin}{7.0\%}
\newcommand{\LongFullBestJoinX}{0.0\%}
% full at best single threshold, complexes aside: identity, split from culture, split from species, joined (all; outside complexes)
\newcommand{\LongFullBestXId}{97.90\%}
\newcommand{\LongFullBestXSplit}{3.3\%}
\newcommand{\LongFullBestXSplitSp}{0.4\%}
\newcommand{\LongFullBestXJoin}{30.0\%}
\newcommand{\LongFullBestXJoinX}{2.9\%}
% full at within-species spread: identity, split from culture, split from species, joined (all; outside complexes)
\newcommand{\LongFullSprId}{98.07\%}
\newcommand{\LongFullSprSplit}{3.9\%}
\newcommand{\LongFullSprSplitSp}{0.6\%}
\newcommand{\LongFullSprJoin}{22.2\%}
\newcommand{\LongFullSprJoinX}{1.9\%}
% full_its at 98.5%: clusters of all copies, of main-type copies, short-only and same-length-only clusters
\newcommand{\ParITSNEAll}{240}
\newcommand{\ParITSNEMain}{220}
\newcommand{\ParITSNEShort}{18}
\newcommand{\ParITSNESame}{2}
\newcommand{\ParITSNEJoin}{2}
% full_its at 97%: clusters of all copies, of main-type copies, short-only and same-length-only clusters
\newcommand{\ParITSNSAll}{110}
\newcommand{\ParITSNSMain}{96}
\newcommand{\ParITSNSShort}{12}
\newcommand{\ParITSNSSame}{2}
\newcommand{\ParITSNSJoin}{2}
% kruger at 98.5%: clusters of all copies, of main-type copies, short-only and same-length-only clusters
\newcommand{\ParKrNEAll}{110}
\newcommand{\ParKrNEMain}{95}
\newcommand{\ParKrNEShort}{13}
\newcommand{\ParKrNESame}{2}
\newcommand{\ParKrNEJoin}{1}
% kruger at 97%: clusters of all copies, of main-type copies, short-only and same-length-only clusters
\newcommand{\ParKrNSAll}{56}
\newcommand{\ParKrNSMain}{46}
\newcommand{\ParKrNSShort}{8}
\newcommand{\ParKrNSSame}{2}
\newcommand{\ParKrNSJoin}{4}
% full at 98.5%: clusters of all copies, of main-type copies, short-only and same-length-only clusters
\newcommand{\ParFullNEAll}{71}
\newcommand{\ParFullNEMain}{58}
\newcommand{\ParFullNEShort}{11}
\newcommand{\ParFullNESame}{2}
\newcommand{\ParFullNEJoin}{3}
% full at 97%: clusters of all copies, of main-type copies, short-only and same-length-only clusters
\newcommand{\ParFullNSAll}{36}
\newcommand{\ParFullNSMain}{30}
\newcommand{\ParFullNSShort}{5}
\newcommand{\ParFullNSSame}{1}
\newcommand{\ParFullNSJoin}{7}
% full ITS at 98.5%: short-only clusters per species that carries the paralog
\newcommand{\ParITSNESpMin}{1}
\newcommand{\ParITSNESpMax}{7}
\newcommand{\ParITSNESpN}{6}
\newcommand{\ParITSNEShared}{0}
% short copies by 5.8S sequence (Glomerales)
\newcommand{\ParShortN}{44}
% SSU flank: identical to a main-type copy
\newcommand{\ParSSUIdentical}{21}
\newcommand{\ParSSUPure}{12}
\newcommand{\ParSSUCrossN}{0}
% full_its: copies in short-only clades of 2+, clades joining taxa, their copies
\newcommand{\ParITSPure}{43}
\newcommand{\ParITSCrossN}{3}
\newcommand{\ParITSCross}{37}
% lsu_flank: copies in short-only clades of 2+, clades joining taxa, their copies
\newcommand{\ParLSUPure}{42}
\newcommand{\ParLSUCrossN}{2}
\newcommand{\ParLSUCross}{10}
% linked: copies in short-only clades of 2+, clades joining taxa, their copies
\newcommand{\ParLinkedPure}{43}
\newcommand{\ParLinkedCrossN}{3}
\newcommand{\ParLinkedCross}{37}
% LSU flank: the R. intraradices/R. prolifer short clade, its support, its sister
\newcommand{\ParLSUIntraProlN}{8}
\newcommand{\ParLSUIntraProlSup}{99/100}
\newcommand{\ParLSUIntraProlSister}{\spn{Sclerocystis sinuosa}}
\newcommand{\ParLSUIntraProlSisterSup}{100/100}
% short-5.8S motif: Rhizophagus short copies (all carry it)
\newcommand{\MotifRhizShort}{41}
% Rhizophagus copies that are not short (none carries it)
\newcommand{\MotifRhizOther}{113}
\newcommand{\MotifSeptoShort}{3}
% Oehlia diaphana copies carrying the motif as their main type
\newcommand{\MotifOehlia}{29}
% retained copies searched
\newcommand{\MotifCopies}{903}
% motif start in the 5.8S of the short copies, and primer starts (fITS7, ITS3)
\newcommand{\MotifStart}{121--122}
\newcommand{\MotifFITSSeven}{84--85}
\newcommand{\MotifITSThree}{31}
% UNITE public release searched for the motif
\newcommand{\UniteSeqs}{2{,}069{,}189}
% Rhizophagus/Rhizoglomus sequences spanning the site (no R. diaphanus)
\newcommand{\UniteRhizCover}{1{,}355}
\newcommand{\UniteRhizShort}{259}
\newcommand{\UniteRhizShortPct}{19\%}
% SHs holding these sequences: only short type, only main type, both
\newcommand{\UniteSHs}{521}
\newcommand{\UniteSHShort}{111}
\newcommand{\UniteSHShortMulti}{32}
\newcommand{\UniteSHMain}{407}
\newcommand{\UniteSHMixed}{3}
% motif in the 5.8S outside Rhizophagus/Rhizoglomus, and in SHs without them
\newcommand{\UniteOther}{984}
\newcommand{\UniteOtherNoRhizSH}{940}
